\chapter{Memory Management and Backend Architecture}
\label{chap:ch3}

\par For most programming languages, memory management is of little concern to the
developer. However, for these two languages, the main similarity is that the
application must be developed with memory usage in mind. Specifically, in C you
must pay close attention to how every bit is allocated, as you are directly
responsible for any small mistake---everything is in your hands. In Rust, on the
other hand, you are constrained by unique development principles designed to keep
everything safe yet equally efficient.

\section{Manual Memory Management in C}
\label{sec:ch3sec1}

In C, memory management requires explicit programmer control. Domain objects can
be allocated on the stack for short-lived values or on the heap using
\texttt{malloc} for longer lifetimes. Each allocation must be paired with a
corresponding \texttt{free} call, and the programmer must track object lifetimes
across function calls.

Local variables on the stack are created when the function starts and released
automatically when it returns. On the heap, the programmer decides both when to
allocate and when to free the memory with \texttt{malloc} and \texttt{free}.
In Gentlix Bank, small helper values such as loop counters stay on the stack,
while long-lived domain objects like \texttt{Account} or \texttt{Transaction}
are allocated on the heap and passed around by pointer.

When dynamic allocation is used, domain entities follow a conventional pattern
with constructor and destructor functions. In our application,
\texttt{account\_create()} internally calls \texttt{malloc}, performs validation,
and frees the allocation immediately if construction fails---the caller never sees
a half-initialised object, but must still call \texttt{account\_free()} on every
successful return:

\clearpage

\begin{lstlisting}[language=C, caption={account.c --- Constructor allocates and
validates. On failure the allocation is freed before returning NULL; on success
the caller owns the pointer and must free it.}]
Account *account_create(UserId user_id, AccountType account_type,
                        Currency currency, int64_t balance_cents,
                        const IBAN *iban, DomainError *err) {
    Account *a = (Account *)malloc(sizeof *a);
    if (!a) { if (err) *err = domain_error_validation("out of memory");
               return NULL; }
    if (!account_init(false, dummy, user_id, account_type,
                      currency, balance_cents, iban, a, err)) {
        free(a);   /* must free before returning on error */
        return NULL;
    }
    return a;      /* caller is now responsible for account_free(a) */
}

void account_free(Account *a) { free(a); }
\end{lstlisting}

This pattern provides flexibility but requires strict discipline---missing a
\texttt{free} call creates memory leaks, while double-free or use-after-free bugs
crash the application. The C approach works well for performance-critical systems
but exposes the entire codebase to memory safety risks.

Another source of bugs in C is the use of \texttt{NULL} pointers. Many functions
return \texttt{NULL} to signal that no value was found or that an allocation
failed. Every call site must test for \texttt{NULL} before dereferencing the
pointer; forgetting this check can easily cause a crash.

A related problem is the dangling pointer. Even after calling \texttt{free}, the
pointer still holds the old memory address. If we use it again by mistake, the
behaviour is undefined---in a banking application this is especially dangerous,
because a wrong read or write can silently corrupt balances or transaction records
without any compiler warning. A missed \texttt{free()} on an error path in the 
transaction service is an example of temporal risk at the application layer; 
Chapter~\ref{chap:ch5} returns to Gentlix with a fuller safety picture.

\newpage
\section{Rust Ownership and Borrowing Model}
\label{sec:ch3sec2}

Rust eliminates manual memory management through its ownership system, enforced
entirely at compile time. Each value has exactly one owner, and when that owner
goes out of scope, Rust automatically invokes destructors via the \texttt{Drop}
trait. This guarantees memory cleanup without garbage collection overhead.

By default, values in Rust live on the stack. A value moves to the heap only
when we wrap it in a smart pointer such as \texttt{Box<T>} or \texttt{Arc<T>}.
In Gentlix Bank, most domain values like a single \texttt{Transaction} start as
stack variables and are moved to the heap only when they need to be shared across
threads or stored behind a repository interface.

Rust's borrowing model allows safe sharing and mutation through references.
Immutable references (\texttt{\&T}) enable read-only access, while mutable
references (\texttt{\&mut T}) guarantee exclusive write access. The compiler
verifies that mutable references never coexist with other references, preventing
data races. The same \texttt{account\_credit}/\texttt{account\_balance}
distinction from C is expressed directly in the method signature---and enforced
by the compiler, not by convention:

\begin{lstlisting}[language=Rust, caption={account.rs --- \texttt{\&self} grants
read-only access; \texttt{\&mut self} grants exclusive mutable access. The
compiler rejects any call that violates these contracts.}]
pub fn balance_cents(&self) -> i64 { self.balance_cents }

pub fn credit(&mut self, amount_cents: i64) -> Result<(), DomainError> {
    if amount_cents <= 0 { return Err(DomainError::validation("...")); }
    self.balance_cents += amount_cents;
    Ok(())
}
\end{lstlisting}

This model results in simpler function contracts: methods take \texttt{\&self}
or \texttt{\&mut self} without exposing allocation details to callers, and the
compiler inserts cleanup code automatically when a value goes out of scope.

Instead of \texttt{NULL} pointers, Rust uses \texttt{Option<T>} to represent
values that may be absent. The compiler forces every call site to handle both
\texttt{Some} and \texttt{None} before the inner value can be used, so a
missing check is a compile-time error, not a runtime crash.

The borrow checker also prevents dangling references entirely. A reference
cannot outlive the value it points to, so storing or returning a reference to
freed memory simply does not compile. The class of bugs described in
Section~\ref{sec:ch3sec1}---dangling pointers, use-after-free,
double-free---does not exist in safe Rust.

\section{Immutability and Shadowing}
\label{sec:ch3sec3}

In C, every local variable is mutable by default. The \texttt{const} keyword
can mark a pointer parameter as read-only, but this protection is limited: it
only restricts access through that specific pointer, and another part of the
program can still cast \texttt{const} away and modify the value. There is no
language mechanism that guarantees a variable will never change after it is
created.

Rust takes the opposite approach. Every binding declared with \texttt{let} is
immutable by default. To allow mutation, the programmer must write
\texttt{let mut} explicitly. This makes the intent clear at the point of
declaration and prevents accidental modifications in long functions.

Rust also supports shadowing: a new \texttt{let} binding can reuse the same
name as an earlier one, hiding the old value from that point forward. This is
useful when a raw value needs to be validated or transformed into a stronger
type. In the service layer of Gentlix Bank, a raw string from the request is
parsed into a validated domain type under the same name:

\begin{lstlisting}[language=Rust, caption={user.rs (service) --- Shadowing
replaces the raw String with a validated Email. After this line, the raw
value is unreachable; only the validated type exists in scope.}]
let email: Email = cmd.email.parse()?;
// cmd.email (raw String) is now shadowed - unreachable from here
\end{lstlisting}

In C, the equivalent transformation requires a separate variable name or
careful discipline to avoid using the raw value after validation. The compiler
offers no help if the programmer accidentally reads the unvalidated string
later in the function.

\section{Lifetime Safety in Practice}
\label{sec:ch3sec4}

The three sections above describe individual mechanisms. Together, they form a
broader difference in how each language treats safety over time.

In C, the programmer is responsible for the full lifetime of every object: when
it is created, when it is valid to read or write, and when it must be freed.
Getting this right is possible, but the compiler does not check it. A pointer
to a freed \texttt{Account} looks identical to a valid one; a missing
\texttt{NULL} check compiles without a warning; a \texttt{const} cast silently
removes a protection. Each of these mistakes produces undefined behaviour that
may only appear in production, under specific conditions.

In Rust, the compiler tracks the lifetime of every value and every reference.
It knows when a value is dropped, whether a reference is still valid, and
whether two references could alias. A program that violates any of these rules
does not compile. This moves an entire class of bugs---use-after-free, double-free,
dangling pointers, data races---from runtime to compile time.

The cost is a steeper learning curve. Beginners often encounter the borrow
checker when they write code that seems correct but violates a lifetime rule.
Over time, working through these errors tends to build better habits around 
ownership and mutation — even if the learning curve feels steep at first. 
In a fintech application like Gentlix Bank, where correctness is critical, 
this trade-off is a strong argument in favour of Rust.

\section*{Summary}

\begin{table}[htbp]
\centering
\begin{tabular}{|l|p{5.5cm}|p{5.5cm}|}
\hline
\textbf{Aspect} & \textbf{C} & \textbf{Rust} \\
\hline
Allocation & Manual \texttt{malloc}/\texttt{free} &
Automatic via ownership \\
\hline
Stack vs.\ heap & Programmer decides for each value
(stack local or heap pointer). Stack values are cheap,
heap values must be freed manually. &
Local variables live on the stack by default; values move
to the heap only through smart pointers such as \texttt{Box}
or \texttt{Arc}. \\
\hline
Sharing & Raw pointers, manual \texttt{NULL} checks &
\texttt{\&T}/\texttt{\&mut T}, compile-time checked \\
\hline
Null / optional values & \texttt{NULL} pointer used to mean
``no value''; every call site must remember to check before
dereference. &
No \texttt{NULL} pointers in safe code; optional values use
\texttt{Option<T>}, and the compiler forces both cases to be
handled. \\
\hline
Cleanup & Explicit \texttt{account\_free()} calls &
Automatic \texttt{Drop} trait \\
\hline
Mutation intent & \texttt{const} pointer convention &
\texttt{\&self}/\texttt{\&mut self} enforced by compiler \\
\hline
Immutability & \texttt{const} qualifiers on pointers and
variables; can be bypassed by casts. &
Default \texttt{let}, explicit \texttt{mut}; borrow checker
enforces read-only or mutable access rules. \\
\hline
Lifetime errors & Runtime crashes or leaks if a pointer is
used after \texttt{free} or if an object is freed too late. &
Compile-time prevention: the compiler tracks lifetimes and
rejects code that could outlive its owners. \\
\hline
Dangling references & Easy to create dangling pointers by
storing or returning a pointer after \texttt{free}; behaviour
is undefined. &
References cannot outlive the data they borrow; creating a
dangling reference is a compile-time error. \\
\hline
\end{tabular}
\caption{Comparison of C and Rust memory management approaches}
\label{tab:memory_comparison}
\end{table}
